\documentclass{article}

\usepackage{amsmath,amssymb}
\usepackage{xurl}
\usepackage[hidelinks]{hyperref}
\setlength{\emergencystretch}{2em}

% Shared commands for notes in docs/paper-gaps/.

\DeclareUnicodeCharacter{2080}{\ifmmode{}_{0}\else\textsubscript{0}\fi}
\DeclareUnicodeCharacter{2081}{\ifmmode{}_{1}\else\textsubscript{1}\fi}
\DeclareUnicodeCharacter{2082}{\ifmmode{}_{2}\else\textsubscript{2}\fi}
\DeclareUnicodeCharacter{2099}{\ifmmode{}_{n}\else\textsubscript{n}\fi}

\newcommand{\Lean}{\textsc{Lean}}
\ExplSyntaxOn
\NewDocumentCommand{\leanid}{m}{
  \ifmmode
    \textsf{\detokenize{#1}}
  \else
    \group_begin:
    \urlstyle{sf}
    \tl_set:Nn \l_tmpa_tl {#1}
    \tl_replace_all:Nnn \l_tmpa_tl {\_} {_}
    \exp_args:NV \path \l_tmpa_tl
    \group_end:
  \fi
}
\ExplSyntaxOff
\providecommand{\tr}{\operatorname{tr}}
\newcommand{\tnleanIssueBase}{https://github.com/LionSR/TNLean/issues/}
\newcommand{\tnleanPrBase}{https://github.com/LionSR/TNLean/pull/}
\newcommand{\tnleanDocBase}{https://sirui-lu.com/TNLean/docs/}
\newcommand{\ghissue}[1]{\href{\tnleanIssueBase#1}{GitHub issue \##1}}
\newcommand{\ghpr}[1]{\href{\tnleanPrBase#1}{PR \##1}}
\newcommand{\doclink}[1]{\href{\tnleanDocBase#1.html}{\path{#1}}}

% Dirac notation and the identity operator, as in blueprint/src/macros/common.tex.
% \providecommand keeps note-local definitions authoritative.
\usepackage{dsfont}
\providecommand{\ket}[1]{|#1\rangle}
\providecommand{\bra}[1]{\langle #1|}
\providecommand{\braket}[2]{\langle #1 | #2 \rangle}
\providecommand{\ketbra}[2]{|#1\rangle\!\langle #2|}
\providecommand{\Id}{\mathds{1}}
\providecommand{\one}{\mathds{1}}
\providecommand{\C}{\mathbb{C}}
\providecommand{\MN}[1]{M_{#1}(\C)}
\providecommand{\MD}{M_D(\C)}

% External referencing for paper sources.  The build helper writes sanitized
% auxiliary files under build/paper-aux/ and excludes duplicated source labels.
\usepackage{xr}

% Import a generated paper auxiliary file with a caller-supplied label prefix.
% The candidate paths support builds from docs/paper-gaps/, the repository root,
% and build/paper-aux/ itself.
\newcommand{\importpaperaux}[2]{%
  \IfFileExists{../../build/paper-aux/#2.aux}{%
    \externaldocument[#1]{../../build/paper-aux/#2}%
  }{%
    \IfFileExists{build/paper-aux/#2.aux}{%
      \externaldocument[#1]{build/paper-aux/#2}%
    }{%
      \IfFileExists{#2.aux}{%
        \externaldocument[#1]{#2}%
      }{}%
    }%
  }%
}

% General external reference macros
\newcommand{\paperref}[2]{\ref{#1-#2}}
\newcommand{\papereqref}[2]{(\ref{#1-#2})}

% CPSV16 source (1606.00608 / MPDO-22-12-17-2.tex) external document and shortcuts
\importpaperaux{cpsv16-}{1606.00608/MPDO-22-12-17-2-tnlean}
\newcommand{\cpsvref}[1]{\paperref{cpsv16}{#1}}
\newcommand{\cpsveqref}[1]{\papereqref{cpsv16}{#1}}

% Verdict marker: \gapnote{<kind>}{<status>}. Kind is the policy
% classification (clarification, local-correction, scope-restriction,
% unfaithful, false-source, open-gap); status is open, wip (elimination
% underway), resolved, or historical. Machine-read by the site index;
% prints a verdict line.
\newcommand{\gapnote}[2]{\par\noindent\textbf{Verdict:} #1 (#2).\par}


\title{The NNCPH ground-state clause in Cirac--Perez-Garcia--Schuch--Verstraete 2016}
\author{}
\date{2026-06-12}

\begin{document}
\maketitle
\gapnote{scope-restriction}{open}

\section{Source statement}

The pure-state main theorem of Ref.~\cite{Cirac2016MPDO_arXiv} states that, for
a tensor $A$ in canonical form generating a family of matrix product states
(MPS), the following conditions are equivalent:\footnote{Local source:
\path{Papers/1606.00608/MPDO-22-12-17-2.tex:534--540}, label
\path{thm:main-MPS}.}
\begin{align}
    \mathrm{RFP}(A),\qquad
    \mathrm{ZCL}(A),\qquad
    \forall N>2,
    \ket{V^{(N)}(A)}
    \text{ is a ground state of a NNCPH}.
    \label{eq:cpsv16_nncph_ground_state_scope_source_equivalence}
\end{align}
Here RFP denotes the renormalization-fixed-point condition, ZCL denotes the
zero-correlation-length condition, and NNCPH denotes a nearest-neighbor
commuting parent Hamiltonian. The formal implications use the same range
$N>2$. This range is independent of the ordered-cycle formula used in the
source argument. At $N=2$, the two periodic length-two windows traverse the
same sites in opposite orders. The sector-graph dimension formula discussed
below handles this convention without extending
(\ref{eq:cpsv16_nncph_ground_state_scope_source_equivalence}) to length two.

The definition preceding the source theorem constructs
\begin{align}
    H_L^{(N)}
        &=\sum_{j=0}^{N-1}\tau_j(P_L^\perp),
    \label{eq:cpsv16_nncph_ground_state_scope_parent_hamiltonian}
\end{align}
where $P_L$ is the orthogonal projector onto the $L$-site MPS subspace and
$\tau_j$ translates an $L$-site operator by $j$ sites. It requires a commuting
parent Hamiltonian to satisfy
\begin{align}
    [\tau_j(P_L),P_L]
        &=0,
        \qquad j=1,\ldots,L-1.
    \label{eq:cpsv16_nncph_ground_state_scope_source_commutation}
\end{align}
A nearest-neighbor parent Hamiltonian has $L=2$.\footnote{Local source:
\path{Papers/1606.00608/MPDO-22-12-17-2.tex:517--524}.}
Define the translated two-site interaction by
\begin{align}
    h_i
        &=\tau_i(P_2^\perp).
    \label{eq:cpsv16_nncph_ground_state_scope_local_interaction}
\end{align}
Under the periodic translation convention, the $L=2$ instance of
(\ref{eq:cpsv16_nncph_ground_state_scope_source_commutation}) gives the
nearest-neighbor commutation relation for the operators $h_i$ after passage
to orthogonal complements. Terms with disjoint two-site supports commute by
locality.

\section{Formal boundary}

For the interaction in
(\ref{eq:cpsv16_nncph_ground_state_scope_local_interaction}), the predicate
\leanid{MPSTensor.IsNNCPH} records
\begin{align}
    h_i h_j
        &=h_j h_i,
        \qquad i,j\in\{0,\ldots,N-1\}.
    \label{eq:cpsv16_nncph_ground_state_scope_formal_commutation}
\end{align}
The predicate \leanid{MPSTensor.IsNNCPHGroundState} also records the
frustration-free ground-vector equations
\begin{align}
    h_i V^{(N)}(A)
        &=0,
        \qquad i\in\{0,\ldots,N-1\}.
    \label{eq:cpsv16_nncph_ground_state_scope_ground_vector}
\end{align}
It therefore captures the annihilation clause of Theorem~3.10(iii) of
Ref.~\cite{Cirac2016MPDO_arXiv} for the canonical projector parent interaction
used in the formal development.

The source definition of a parent Hamiltonian also requires, for $N>L$, that
the periodic MPS vectors span the entire ground-state subspace. The formal
development names this condition
\leanid{MPSTensor.HasParentHamiltonianGroundSpaceSpanning}; it is not part of
\leanid{MPSTensor.IsNNCPHGroundState}. The fixed-chain predicate
\leanid{MPSTensor.HasNNCPHGroundSpace} combines
(\ref{eq:cpsv16_nncph_ground_state_scope_formal_commutation}),
(\ref{eq:cpsv16_nncph_ground_state_scope_ground_vector}), and
\begin{align}
    \ker H_2^{(N)}(B)
        &=
        \operatorname{span}
        \{V^{(N)}(A_j):j=1,\ldots,g\}.
    \label{eq:cpsv16_nncph_ground_state_scope_ground_space_spanning}
\end{align}
The source theorem requires the all-chain predicate
\leanid{MPSTensor.HasNNCPHGroundSpaces}, which quantifies this fixed-chain
condition over $N>2$. This predicate names a condition; it does not establish
(\ref{eq:cpsv16_nncph_ground_state_scope_ground_space_spanning}) for a
concrete canonical-form tensor. The declaration
\leanid{MPSTensor.hasNNCPHGroundSpaces\_iff\_forall\_isNNCPHGroundState\_and\_groundSpaceSpanning}
formalizes the separation: the all-chain source condition is exactly the
all-chain ground-vector condition together with the nearest-neighbor
ground-space spanning clause.

The easy inclusion in
(\ref{eq:cpsv16_nncph_ground_state_scope_ground_space_spanning}) is proved.
If every component vector $V^{(N)}(A_j)$ is annihilated by the parent
Hamiltonian, or more strongly by every local term, then
\begin{align}
    \operatorname{span}
    \{V^{(N)}(A_j):j=1,\ldots,g\}
        &\subseteq
        \ker H_L^{(N)}(B).
    \label{eq:cpsv16_nncph_ground_state_scope_easy_kernel_inclusion}
\end{align}
\footnote{Formal declarations:
\leanid{MPSTensor.bntMPSVectorSpan\_le\_ker\_parentHamiltonian\_of\_forall\_annihilates}
and
\leanid{MPSTensor.bntMPSVectorSpan\_le\_ker\_parentHamiltonian\_of\_forall\_frustrationFree}.}
For a block-diagonal tensor with nonzero weights,
\leanid{MPSTensor.hasParentHamiltonianGroundSpaceSpanning\_toTensorFromBlocks\_iff\_ker\_le\_bntMPSVectorSpan}
reduces the spanning equality to the reverse containment, using
(\ref{eq:cpsv16_nncph_ground_state_scope_easy_kernel_inclusion}). Its
nearest-neighbor specialization,
\leanid{MPSTensor.hasNNCPHGroundSpaces\_toTensorFromBlocks\_iff\_forall\_isNNCPH\_and\_ker\_le\_bntMPSVectorSpan},
states the all-chain condition as translated two-site commutation together
with the same reverse containment for every $N>2$.

For the direct sum containing one copy of every distinct BNT basis tensor,
transfer idempotence of the whole tensor implies the reverse containment. The
resulting theorem is
\leanid{MPSTensor.IsBNTCanonicalForm.rfp\_hasParentHamiltonianGroundSpaceSpanning\_basisDirectSum}.

The commutation clause is also proved for this multiplicity-one direct sum.
The residual matrix entries form an orthonormal family on the disjoint union
of the within-sector matrix spaces. The equal-sector virtual bond projector
is the sum of the sector bond projectors, extended by zero between unequal
sectors. Its adjacent placements commute sectorwise, and every mixed product
vanishes by orthogonality of the canonical summand inclusions. Transport by
the common physical isometry identifies this virtual projector with the
canonical projector onto the direct sum's two-site space. Complementing and
transporting the resulting three-site relation around the periodic chain gives
\leanid{MPSTensor.IsBNTCanonicalForm.rfp\_implies\_nncph\_basisDirectSum}.

Combining commutation with ground-space spanning gives
\leanid{MPSTensor.IsBNTCanonicalForm.rfp\_implies\_hasNNCPHGroundSpaces\_basisDirectSum}.
This implication assumes only BNT canonical form and transfer idempotence of
the whole multiplicity-one direct sum. It uses no comparison, crossing,
repeated-copy, raw-weight, or externally supplied commutation hypothesis.

This result does not prove the full source theorem. Together with the reverse
implication and the corrected positive-gap ZCL equivalence, it proves the
three-way equivalence for the multiplicity-one, unit-weight basis direct sum.
Repeated copies, arbitrary raw sector weights, adjacent complementary gaps,
and the unrestricted canonical-form statement remain excluded.

\section{Forward implication and Appendix B structure}

The three-site $AX/XB$ support maps used in the forward
RFP-to-NNCPH implication are represented by
\leanid{MPSTensor.HasOverlappingTwoSiteCommutation}. Appendix~B of
Ref.~\cite{Cirac2016MPDO_arXiv} supplies the basic-vector form. The local
projectors in Definition~D.2 of Ref.~\cite{Cirac2016MPDO_arXiv} belong to a
separate Appendix~D argument and are not invoked by the one-sentence proof of
RFP $\Rightarrow$ NNCPH at source lines 1305--1307.

The basic-vector coefficients define a physical isometry $U$ by
\begin{align}
    (Uv)_i
        &=\sum_{\alpha,\beta}
        U^i_{\alpha,\beta}v_{\alpha,\beta},
    \label{eq:cpsv16_nncph_ground_state_scope_physical_isometry_action}\\
    U^*U
        &=1.
    \label{eq:cpsv16_nncph_ground_state_scope_physical_isometry_left_inverse}
\end{align}
These identities are recorded by
\leanid{MPSTensor.AppendixBStructuralData.physicalIsometry} and
\leanid{MPSTensor.AppendixBStructuralData.physicalIsometryLeftInverse\_comp}.
For every length $L$, the tensor power and its conjugate coefficient map obey
\begin{align}
    U^{*\otimes L}U^{\otimes L}
        &=1.
    \label{eq:cpsv16_nncph_ground_state_scope_tensor_power_left_inverse}
\end{align}
The corresponding declaration is
\leanid{MPSTensor.AppendixBStructuralData.physicalIsometryTensorPowerLeftInverse\_comp}.

Define the two-site bond insertion $I_\varphi$ by
\begin{align}
    (I_\varphi v)_{(a,b),(b',c)}
        &=\delta_{b,b'}\lambda_b v_{a,c}.
    \label{eq:cpsv16_nncph_ground_state_scope_bond_insertion}
\end{align}
Transport through $U^{\otimes2}$ gives
\begin{align}
    \operatorname{ran}(U^{\otimes2}I_\varphi)
        &=G_2(A_{\mathrm{core}}).
    \label{eq:cpsv16_nncph_ground_state_scope_two_site_range}
\end{align}
This range equality is recorded by
\leanid{MPSTensor.AppendixBStructuralData.twoSiteBasicEmbedding\_range}. The
orthogonal projector onto this range is
\leanid{MPSTensor.AppendixBStructuralData.twoSiteBasicSupportProjection}; it
satisfies
\begin{align}
    P_2
        &=1-q_2(A_{\mathrm{core}}).
    \label{eq:cpsv16_nncph_ground_state_scope_support_projection_complement}
\end{align}
The declaration
\leanid{MPSTensor.AppendixBStructuralData.twoSiteBasicSupportProjection\_eq\_complement}
records this identity.

The coefficient-space maps
\leanid{MPSTensor.AppendixBStructuralData.appendixBQAXOnCoeffSpace} and
\leanid{MPSTensor.AppendixBStructuralData.appendixBQXBOnCoeffSpace} are
identified with the canonical two-site parent interaction $q_2(A)$ before
placement on the $AX$ and $XB$ faces. Denote these maps by
$\widehat Q_{AX}$ and $\widehat Q_{XB}$. Their lifts act on the three-site
coefficient space; they have not yet been identified with the source
projectors on $\mathcal H_A\otimes\mathcal H_X$ and
$\mathcal H_X\otimes\mathcal H_B$.

The one-bond virtual projector and its adjacent placements are constructed in
\leanid{MPSTensor.AppendixBStructuralData.virtualBondProjection},
\leanid{MPSTensor.AppendixBStructuralData.leftVirtualBondProjection}, and
\leanid{MPSTensor.AppendixBStructuralData.rightVirtualBondProjection}.
Transporting their commutation through $U^{\otimes3}$ gives commutation of the
adjacent physical support projectors, formalized by
\leanid{MPSTensor.AppendixBStructuralData.twoSiteBasicSupportProjection\_commute\_lifts},
and hence the unconditional coefficient-space relation
\leanid{MPSTensor.AppendixBStructuralData.hasOverlappingTwoSiteCommutation}.

The cyclic three-site restriction transports this commutator to every adjacent
pair on every chain with $N>2$. The relevant declarations are
\leanid{MPSTensor.localTerm\_adjacent\_twoSite\_commute\_of\_threeSite\_zero\_one\_commute}
and
\leanid{MPSTensor.AppendixBStructuralData.adjacent\_twoSite\_localTerms\_commute}.
The resulting commuting projector family is
\leanid{MPSTensor.AppendixBStructuralData.hasProductPairLocalProjectors}.

The coefficient-space analogues of the Appendix~D algebraic and
kernel-intersection equations are also proved locally. The declaration
\leanid{MPSTensor.AppendixBStructuralData.hasAppendixD2ParentCommutingHamiltonian}
identifies $\widehat K_{AXB}=G_3(A)$ with the intersection of the kernels of
the lifts of $\widehat Q_{AX}$ and $\widehat Q_{XB}$. Its RFP corollary derives
injectivity from the normal left-canonical RFP hypotheses. The concrete
canonical parent interactions are orthogonal by construction, but this
declaration does not assert self-adjointness. Constructing the source
projectors and identifying their transported actions with the hatted maps
remains open. This local coefficient-space result is not used in the chain
commutator proof.

The declarations
\leanid{MPSTensor.rfp\_implies\_nncph\_of\_leftCanonical} and
\leanid{MPSTensor.rfp\_implies\_nncph\_ground\_state\_of\_leftCanonical}
therefore prove the forward commutation and zero-energy statements for normal
left-canonical RFP tensors. For a CPSV normal tensor, spectral-radius-one
normalization supplies a pure trace-preserving gauge. Gauge invariance of the
physical RFP relation and of the canonical parent interactions transports the
result to \leanid{MPSTensor.rfp\_implies\_nncph} and
\leanid{MPSTensor.rfp\_implies\_nncph\_ground\_state}. These declarations prove
commutation and zero energy for a single CPSV normal tensor without an
external axiom.

For the same tensor, transfer idempotence and normality imply one-site
injectivity. The injective periodic-chain theorem identifies the kernel of the
two-site parent Hamiltonian with the line spanned by the MPS vector for every
$N>2$. This result is
\leanid{MPSTensor.rfp\_hasParentHamiltonianGroundSpaceSpanning\_single}.
Combining this equality with commutation and zero energy gives
\leanid{MPSTensor.rfp\_implies\_hasNNCPHGroundSpaces\_single}. Thus the complete
forward NNCPH ground-space condition is proved for one normal sector.

\section{Multiplicity-one basis direct sums}

Let a BNT canonical form contain several distinct basis sectors. Transfer
idempotence of their direct sum gives residual tensors whose one-site entry
vectors have identity Gram matrix. The simultaneous one-site word tuples
therefore span the product of all block matrix algebras. This span gives a
common one-letter expression for every block identity and propagates to every
positive word length, replacing the right-canonical hypothesis in the local
restriction-intersection calculation.

A global one-site change of cut then forces the block boundary matrices to
commute with their tensors. Consequently, the periodic nearest-neighbor chain
space splits as the sum of the single-block chain spaces. The single-block
injective periodic-chain theorem then proves
\leanid{MPSTensor.IsBNTCanonicalForm.rfp\_hasParentHamiltonianGroundSpaceSpanning\_basisDirectSum}.
The independent residual-family bond-projector argument proves
\leanid{MPSTensor.IsBNTCanonicalForm.rfp\_implies\_nncph\_basisDirectSum}.
Their conjunction is
\leanid{MPSTensor.IsBNTCanonicalForm.rfp\_implies\_hasNNCPHGroundSpaces\_basisDirectSum}.
This is exactly the multiplicity-one, distinct-sector forward case; it does
not include repeated copies or arbitrary raw weights.

The reverse theorem \leanid{MPSTensor.nncph\_implies\_rfp} assumes both the
explicit BNT relation \leanid{MPSTensor.IsBNT} and the all-chain condition
\leanid{MPSTensor.HasNNCPHGroundSpaces}. It also assumes that the represented
tensor is normal and left-canonical. These hypotheses select a normal
representative and fix its normalization; they do not reproduce the full
canonical-form hypothesis of Theorem~3.10 of
Ref.~\cite{Cirac2016MPDO_arXiv}.

The reverse implication is also proved for the multiplicity-one, unit-weight
direct sum of all distinct BNT basis sectors. The declaration
\leanid{MPSTensor.IsBNTCanonicalForm.isTransferIdempotent\_basisDirectSum\_of\_hasNNCPHGroundSpaces}
matches every BNT basis tensor with a distinct normalized loop tensor from the
finite-degeneracy classification used in Ref.~\cite{Cirac2016MPDO_arXiv}. The
disjoint loop sectors have zero mixed transfer maps, and this local
orthogonality is preserved by the independent virtual gauge-phase relations.
The direct-sum transfer criterion then gives transfer idempotence. This result
still excludes repeated copies and arbitrary raw weights.

\section{Sector-graph reduction}

Let $P$ be a sector decomposition with representatives
$A_1,\ldots,A_g$ and total tensor $B$. If the representative MPS vectors are
eventually linearly independent, the all-chain condition gives
\begin{align}
    \dim_{\C}\ker H_2^{(N)}(B)
        &=g
        \qquad\text{for all sufficiently large }N.
    \label{eq:cpsv16_nncph_ground_state_scope_eventual_ground_dimension}
\end{align}
This result is
\leanid{MPSTensor.HasBNTSectorData.eventually\_parentHamiltonian\_groundSpace\_finrank\_eq}.
Its proof uses only ground-space spanning and the dimension of the span of a
finite linearly independent family; it does not use the commuting-interaction
classification.

The canonical two-site parent interaction
\leanid{MPSTensor.twoSiteParentInteractionMatrix} is Hermitian because it is
the transport of an orthogonal projector. When the nearest-neighbor parent
terms commute on three sites, their common left-associated coordinates are
the two natural overlapping lifts of this matrix. Hence
\leanid{MPSTensor.IsNNCPH.twoSiteParentInteractionMatrix\_isHermitian\_and\_overlappingLifts\_commute}
supplies the hypotheses of the spatial decomposition used in
Ref.~\cite{Cirac2016MPDO_arXiv}, without any canonical-form or
ground-space-spanning assumption. Its consequence
\leanid{MPSTensor.IsNNCPH.exists\_unitary\_blockActions\_of\_twoSiteParentInteractionMatrix}
decomposes the middle physical space into a finite direct sum of left and
right tensor factors on which the adjacent interactions act separately.

The symmetric two-sided form is also internal. The complementary two-site
interaction is a positive semidefinite orthogonal projector, and its
overlapping lifts commute under the same three-site hypothesis. The theorem
\leanid{Matrix.exists\_positive\_etaPair\_decomposition\_of\_overlappingLifts\_commute}
combines the one-sided block actions into
\begin{align}
    (U\otimes U)^*(\one-Q_A)(U\otimes U)
        &=
        \bigoplus_{q,h}
        \one_{q,l}\otimes\eta_{q,h}\otimes\one_{h,r}.
    \label{eq:cpsv16_nncph_ground_state_scope_symmetric_block_decomposition}
\end{align}
Idempotence of the left-hand side makes every $\eta_{q,h}$ idempotent;
positivity then makes it an orthogonal projector. Consequently,
\leanid{MPSTensor.IsNNCPH.nonempty\_beigiSectorGraphData} constructs a finite
sector graph directly from the three-site commuting-parent hypothesis, and
\leanid{MPSTensor.IsNNCPH.beigiSectorGraphData} selects such a graph. This
construction uses no canonical-form, all-chain ground-space, or
eventual-dimension hypothesis.

For a finite sector graph,
\leanid{MPSTensor.BeigiSectorGraphData.parentHamiltonianGroundSpaceES\_finrank\_eq}
proves the ordered-cycle ground-space dimension formula for every $N\geq2$.
Its period-two specialization,
\leanid{MPSTensor.BeigiSectorGraphData.parentHamiltonianGroundSpaceES\_finrank\_eq\_two},
retains both orientations: when $a\ne b$, the ordered cycles $(a,b)$ and
$(b,a)$ are distinct. This is the convention used in Section~III,
Equations~(2)--(4), of the finite-degeneracy argument cited by
Ref.~\cite{Cirac2016MPDO_arXiv}. These sector-graph theorems assume no
chain-level ground-space equality, eventual degeneracy, or canonical-form
datum.

At $N=1$, the cited two-site interaction has no specified action. The
finite-degeneracy argument used by Ref.~\cite{Cirac2016MPDO_arXiv} defines
$P_{j,j+1}$ on sites $j,j+1$ and separates the actions of $P_{j-1,j}$ and
$P_{j,j+1}$ on the middle site. At $N=1$, all these site labels coincide, and
the source gives no rule for converting the two-site operator into a one-site
operator. Theorem~3.10(iii) of Ref.~\cite{Cirac2016MPDO_arXiv} likewise
concerns only $N>2$.

The project uses the local convention that an $L$-site interaction is zero
when $L>N$. Thus
\leanid{MPSTensor.parentHamiltonian\_two\_one\_eq\_zero} proves
$H_2^{(1)}=0$, and
\leanid{MPSTensor.parentHamiltonianGroundSpaceES\_two\_one\_eq\_top} proves
\begin{align}
    \ker H_2^{(1)}
        &=\C^d.
    \label{eq:cpsv16_nncph_ground_state_scope_one_site_kernel}
\end{align}
This convention does not extend the ordered-cycle formula. An ordered
length-one cycle is a self-loop, so the graph expression would be
\begin{align}
    \sum_a\dim\operatorname{ran}P_{a,a}.
    \label{eq:cpsv16_nncph_ground_state_scope_self_loop_dimension}
\end{align}
The two-site block decomposition does not equate
(\ref{eq:cpsv16_nncph_ground_state_scope_self_loop_dimension}) with $d$.

For example, consider the tensor
\begin{align}
    A^0
        &=
        \begin{pmatrix}
            0&1\\
            0&0
        \end{pmatrix},
    \label{eq:cpsv16_nncph_ground_state_scope_counterexample_A_zero}\\
    A^1
        &=
        \begin{pmatrix}
            0&0\\
            1&0
        \end{pmatrix}.
    \label{eq:cpsv16_nncph_ground_state_scope_counterexample_A_one}
\end{align}
Its two-site ground space is spanned by $\ket{01}$ and $\ket{10}$. The
corresponding one-dimensional sectors have only the edges $0\to1$ and
$1\to0$. Hence
(\ref{eq:cpsv16_nncph_ground_state_scope_self_loop_dimension}) is zero, while
the project-local one-site ground space in
(\ref{eq:cpsv16_nncph_ground_state_scope_one_site_kernel}) has dimension
$d=2$. The formal development therefore treats the project-local result
separately and proves the ordered-cycle theorem only in its source-defined
range $N\geq2$, including both orientations at $N=2$. The resolved source
boundary is recorded in \ghissue{4400}: no cited source defines a two-site
operator $H_1$ or a coincident periodic-edge convention at $N=1$.
\ghissue{4655} records the separate project-local boundary.

The remaining finite graph comparison is also proved without an axiom. For a
finite directed graph with nonnegative integral edge weights, the declarations
\leanid{FiniteWeightedDigraph.cyclicEdgeWeight\_eq\_one},
\leanid{FiniteWeightedDigraph.cycle\_eq\_constant}, and
\leanid{FiniteWeightedDigraph.cycleWeightSum\_eq\_card\_loop}
show that eventual constancy of the ordered-cycle sum forces every cyclic edge
weight to equal one, every directed cycle to be a repeated loop, and every
positive-length cycle sum to equal the number of positive loops. The proof
uses the two repetition maps from consecutive sufficiently large lengths into
their common product, as in the source argument.

Applied to the sector datum,
\leanid{MPSTensor.BeigiSectorGraphData.parentHamiltonianGroundSpaceES\_finrank\_eq\_card\_loop}
identifies every parent-ground-space dimension for $N>2$ with the loop number
when these dimensions are eventually constant. Combining this result with
(\ref{eq:cpsv16_nncph_ground_state_scope_eventual_ground_dimension}) gives
\leanid{MPSTensor.BeigiSectorGraphData.parentHamiltonianGroundSpaceES\_finrank\_eq\_card\_loop\_of\_hasBNTSectorData}.
This conclusion does not identify the loop states with the original BNT
representatives.

\section{Loop tensors and ground-space spanning}

For each positive loop, the product-of-pairs vector is constructed explicitly.
The declaration \leanid{MPSTensor.BeigiSectorGraphData.loopTensor} uses the
loop sector's left factor as its virtual space and denotes the chosen nonzero
loop-edge vector by $\varphi_a$. At site $n$, let $r_n$ and $l_n$ denote the
right and left sector factors. The cyclic bond pairs $r_n$ with $l_{n+1}$. For
every positive length,
\leanid{MPSTensor.BeigiSectorGraphData.mpv\_loopTensor} identifies the periodic
MPS vector with the sitewise-unitary image of
\begin{align}
    \prod_n\varphi_a(r_n,l_{n+1}).
    \label{eq:cpsv16_nncph_ground_state_scope_loop_product_state}
\end{align}
This is the product-of-pairs construction in Section~IV, Equation~(13), of the
finite-degeneracy argument cited by Ref.~\cite{Cirac2016MPDO_arXiv}. For a
positive loop $a\to a$, denote the resulting length-$N$ physical state by
$\Psi_a^{(N)}$.

The ambient left factor need not be the minimal virtual space. The construction
requires only that the selected bond vector be nonzero, not that its
coefficient matrix have full rank. Factor its coefficient matrix through its
positive Schmidt rank $r_a$ as
\begin{align}
    \Phi_a
        &=X_aY_a.
    \label{eq:cpsv16_nncph_ground_state_scope_schmidt_rank_factorization}
\end{align}
The declarations
\leanid{MPSTensor.BeigiSectorGraphData.minimalLoopCoordinateTensor} and
\leanid{MPSTensor.BeigiSectorGraphData.minimalLoopTensor} restrict the virtual
space to $\C^{r_a}$. Their tensor letters are
$(Y_a)_{\alpha,l}(X_a)_{r,\beta}$, and
\leanid{MPSTensor.BeigiSectorGraphData.mpv\_minimalLoopTensor} proves that the
restriction leaves every positive-length periodic vector equal to
(\ref{eq:cpsv16_nncph_ground_state_scope_loop_product_state}).

Exactness of
(\ref{eq:cpsv16_nncph_ground_state_scope_schmidt_rank_factorization}) implies
that the columns of $Y_a$ and the rows of $X_a$ span $\C^{r_a}$. Their outer
products are the minimal tensor's letters. Hence
\leanid{MPSTensor.BeigiSectorGraphData.minimalLoopCoordinateTensor\_isInjective}
proves one-site injectivity. The one-site unitary preserves this property,
giving
\leanid{MPSTensor.BeigiSectorGraphData.minimalLoopTensor\_isInjective} and
\leanid{MPSTensor.BeigiSectorGraphData.minimalLoopTensor\_isNormal}.

Here $\lVert\varphi_a\rVert_2$ is the Euclidean norm of the bond vector,
equivalently the Frobenius norm of its coefficient matrix, rather than the
supremum norm on the ordinary coordinate-function space. Define
\begin{align}
    P
        &=Y_aY_a^*,
    \label{eq:cpsv16_nncph_ground_state_scope_transfer_rank_one_P}\\
    Q
        &=X_a^*X_a.
    \label{eq:cpsv16_nncph_ground_state_scope_transfer_rank_one_Q}
\end{align}
The transfer map of the minimal coordinate tensor then satisfies
\begin{align}
    \mathcal E_a(Z)
        &=\tr(QZ)P,
    \label{eq:cpsv16_nncph_ground_state_scope_transfer_rank_one}\\
    \mathcal E_a^2
        &=\lVert\varphi_a\rVert_2^2\mathcal E_a.
    \label{eq:cpsv16_nncph_ground_state_scope_transfer_quasi_idempotence}
\end{align}
The second identity is formalized by
\leanid{MPSTensor.BeigiSectorGraphData.transferMap\_minimalLoopCoordinateTensor\_comp\_self}.

The declaration \leanid{MPSTensor.BeigiSectorGraphData.loopBondVector}
chooses a nonzero vector but does not normalize it. Transfer idempotence of
the raw tensor therefore requires $\lVert\varphi_a\rVert_2=1$. The separate
declaration
\leanid{MPSTensor.BeigiSectorGraphData.normalizedMinimalLoopTensor} defines
$\lVert\varphi_a\rVert_2^{-1}\widetilde C_a$ without changing the raw tensor or
its exact state formulas. Its transfer map is idempotent by
\leanid{MPSTensor.BeigiSectorGraphData.normalizedMinimalLoopTensor\_isTransferIdempotent},
and \leanid{MPSTensor.BeigiSectorGraphData.mpv\_normalizedMinimalLoopTensor}
identifies its length-$N$ MPS vector as
$\lVert\varphi_a\rVert_2^{-N}$ times the original loop state. No identification
with the original basis of normal tensors is asserted at this stage. The
spanning result below may therefore be formulated using the periodic vectors
of the minimal Schmidt-support tensors.

For every $N\geq2$,
\leanid{MPSTensor.BeigiSectorGraphData.loopProductState\_mem\_ker\_parentHamiltonian}
proves that the loop product vector belongs to the kernel of the length-two
periodic parent Hamiltonian. Its Euclidean-space form is
\leanid{MPSTensor.BeigiSectorGraphData.loopProductState\_mem\_parentHamiltonianGroundSpaceES}.
The proof uses the local edge-ground-space condition and its positive-loop
specialization in the finite-degeneracy construction cited by
Ref.~\cite{Cirac2016MPDO_arXiv}. The condition $N\geq2$ is the source-defined
range of the two-site periodic Hamiltonian. As recorded by \ghissue{4400}, no
source-defined coincident-edge convention extends this statement to $N=1$.

At every positive length, the loop vectors are nonzero by
\leanid{MPSTensor.BeigiSectorGraphData.loopProductStateES\_ne\_zero}. Distinct
loop vectors are orthogonal by
\leanid{MPSTensor.BeigiSectorGraphData.loopProductStateES\_inner\_eq\_zero\_of\_ne},
and hence linearly independent by
\leanid{MPSTensor.BeigiSectorGraphData.loopProductStateES\_linearIndependent}.
If the parent-ground-space dimension is eventually constant, membership and
the loop-count formula give, for every $N>2$,
\begin{align}
    \operatorname{span}
    \{\Psi_a^{(N)}:a\to a\text{ is a positive loop}\}
        &=
        \ker H_2^{(N)}(A).
    \label{eq:cpsv16_nncph_ground_state_scope_loop_ground_space_spanning}
\end{align}
This equality is
\leanid{MPSTensor.BeigiSectorGraphData.span\_loopProductStateES\_eq\_parentHamiltonianGroundSpaceES}.
Its BNT and all-chain specialization is
\leanid{MPSTensor.BeigiSectorGraphData.span\_loopProductStateES\_eq\_parentHamiltonianGroundSpaceES\_of\_hasBNTSectorData}.
These results do not by themselves identify the loop states with the original
BNT representatives.

For a normal left-canonical representative, the final comparison is proved.
Algebraic normality implies spectral normality, while the normalized minimal
loop tensors are normal and transfer-idempotent. The all-chain ground-space
equality places the original tensor's MPS vector eventually in the span of the
loop-tensor vectors. The normal-tensor overlap dichotomy and asymptotic
orthogonality of distinct loops then imply that the original tensor agrees
with one loop tensor up to a unit phase and an invertible virtual change of
basis. Transfer idempotence is invariant under both transformations. This
proves \leanid{MPSTensor.nncph\_implies\_rfp}.

\section{Comparison scope and eventual span equality}

The single-normal theorem is a scope restriction relative to Theorem~3.10 of
Ref.~\cite{Cirac2016MPDO_arXiv}, source lines 534--541 and its proof at line
1307. The formal theorem explicitly assumes that the representative is normal
and left-canonical, together with an explicit BNT relation and the all-chain
ground-space condition. It does not state the unrestricted three-way
equivalence for an arbitrary canonical-form tensor. In the source canonical
decomposition, normality selects the one-sector case $g=1$.

The basis-direct-sum theorem establishes the $g>1$ reverse implication for the
multiplicity-one, unit-weight representative by using the finite-degeneracy
classification. Neither theorem includes repeated copies or arbitrary raw
weights.

The parent-ground-space equality in Ref.~\cite{Cirac2016MPDO_arXiv} applies
for $N>2$, whereas the earlier permutation-rigidity lemmas assumed equality
of the two MPS-vector spans at every length. Their proofs use only eventual
span equality: one chooses a sufficiently large length at which both families
are linearly independent, and every subsequent overlap comparison is
asymptotic. The comparison is therefore now stated under eventual span
equality. The source supplies this hypothesis on the concrete tail $N>2$ at
lines 522--525 and 1305--1307. The general comparison lemmas assert only the
existence of a finite tail because the chosen length must also exceed both
linear-independence thresholds. They do not assert a common numerical
threshold of $2$.

This replacement of an unnecessary all-length hypothesis by the finite-tail
hypothesis used in the proof is a local correction. Once the normal
representatives are matched, their gauge-phase relation determines their
vectors at every positive length.

The auxiliary declaration
\leanid{MPSTensor.BeigiSectorGraphData.mem\_ker\_parentHamiltonian\_of\_groundBondProduct\_mulVec\_eq\_self}
assumes sector data, which supplies commutativity of all translated two-site
parent interactions. It is therefore a restricted fixed-vector criterion,
not an unrestricted theorem for arbitrary parent interactions. Any independent
use should replace the sector-data assumption by explicit pairwise
commutativity of the translated interactions.

\section{Completed representative equivalence and residual scope}

For the multiplicity-one, unit-weight direct sum
$A_{\mathrm{basis}}$ of the distinct BNT basis tensors, the corrected
three-way equivalence is
\begin{align}
    \mathrm{RFP}(A_{\mathrm{basis}})
        &\Longleftrightarrow
        \text{positive-gap BNT ZCL}(A_{\mathrm{basis}})
    \label{eq:cpsv16_nncph_ground_state_scope_representative_rfp_zcl}\\
        &\Longleftrightarrow
        \mathrm{HasNNCPHGroundSpaces}(A_{\mathrm{basis}}).
    \label{eq:cpsv16_nncph_ground_state_scope_representative_nncph}
\end{align}
The reverse implication is proved both for a normal left-canonical
representative and for the multiplicity-one, unit-weight direct sum of all
distinct BNT basis sectors. In the direct-sum theorem, the comparison with the
locally orthogonal product-of-pairs loop states is familywise and requires no
axiom. The pairwise equivalences are formalized by
\leanid{MPSTensor.IsBNTCanonicalForm.isTransferIdempotent\_basisDirectSum\_iff\_hasNNCPHGroundSpaces}
and
\leanid{MPSTensor.IsBNTCanonicalForm.isPositiveGapBNTZCL\_basisDirectSum\_iff\_hasNNCPHGroundSpaces}.
\footnote{Tracked by \ghissue{2633}.}

The proved reverse implication therefore covers the corrected basis-direct-sum
representative under the source all-chain ground-space condition. Repeated
copies and arbitrary raw weights remain excluded by the known counterexample;
they are not pending extensions of this representative theorem. The printed
unrestricted equivalence remains unformalized.

For the forward commutation implication, the formal development contains the
source-unit Appendix~B structural datum, its physical isometric embedding, all
tensor powers with explicit left inverses, the two-site bond insertion, the
identification of its physical image with $G_2(A_{\mathrm{core}})$, the
orthogonal support projector, the three-site $AX/XB$ support maps, the
coefficient form of $U^{\otimes L}\varphi_j^{\otimes L}$, the adjacent virtual
and physical projectors, their three-site commutator, transport to every
periodic chain with $N>2$, and the pure trace-preserving gauge passage for a
CPSV normal tensor. The even-chain physical-pair factorization remains a
separate conditional construction and is unnecessary for the proved chain
commutator.

Extension from one normal tensor to the full repeated-sector canonical form
is obstructed as recorded in
\path{docs/paper-gaps/cpsv16_rfp_isometry_scope.tex}. The single-normal-sector
ground-space spanning clause and the multiplicity-one distinct-sector
extension are proved. Repeated copies and arbitrary raw weights lie outside
the corrected theorem.

In the reverse direction, the formal development proves eventual
ground-space-dimension stability, the symmetric sector-graph construction from
three-site commuting parent terms, the ordered-cycle dimension formula, the
reduction to one-dimensional loops, and the exact MPS realization and
ground-space membership of every positive loop state. Under the
eventual-dimension hypothesis supplied by BNT eventual linear independence and
the source all-chain ground-space condition, the loop vectors span the parent
ground space for every $N>2$. Their comparison with the original tensor is
proved for one normal left-canonical representative and familywise for the
multiplicity-one, unit-weight BNT basis direct sum; the corresponding
transfer-idempotence results are also proved.

The two biconditionals in
(\ref{eq:cpsv16_nncph_ground_state_scope_representative_rfp_zcl}) and
(\ref{eq:cpsv16_nncph_ground_state_scope_representative_nncph}) assemble these
representative-level implications. The residual gap is the printed
unrestricted statement. Its repeated-copy, raw-weight, and adjacent-gap cases
are excluded from the corrected theorem.

\bibliographystyle{alpha}
\bibliography{references}

\end{document}
